% Front-matter entry check: questions and answers stay together before Part 1.
\Needspace{.3\textheight}
\section*{开读前：五题入门自检}
\phantomsection\label{app:math-entry}
\begingroup\small
\setlength{\parskip}{3pt}
花五分钟独立完成，不用记新术语。五题都能说明理由，就可以直接开读；
卡在哪项，先复习\hyperref[app:math-tools]{附录~\ref*{app:math-tools} 的数学速查（第~\pageref*{app:math-tools}~页）}中对应的运算。
第 5 题只要求一元导数，链式法则会在正文继续练习。
\begin{enumerate}[itemsep=1pt,topsep=2pt]
\item 解 $2w+1=5$。
\item 计算 $(-2)^2+1$，并说明它与 $-2^2+1$ 的区别。
\item 对 $f(x)=2x+1$，求 $f(-1),f(0),f(1)$。
\item 计算 $\ln(e^2)$，并求满足 $e^u=1$ 的 $u$。
\item 对 $g(w)=(w-2)^2$，求 $g'(w)$ 与 $g'(1)$；应增大还是减小 $w$ 才先靠近最低点？
\end{enumerate}

\subsection*{核对结果与补读建议}
\phantomsection\label{answer:math-entry}
做完再看。\textbf{第一步：}先代入或整理等式，再判断运算顺序。\par\nobreak
\textbf{结果：}1. $w=2$；2. $5$ 与 $-3$，先算幂再应用前面的负号；
3. $(-1,1,3)$；4. $2$ 与 $u=0$；
5. $g'(w)=2(w-2)$，$g'(1)=-2$，导数为负，应先增大 $w$。\par
\textbf{补读：}前三题卡住，先复习附录~\ref*{app:math-tools} 的符号表；第 4 题复习指数与自然对数；第 5 题复习导数与链式法则。
\endgroup
